%% kapa: exam=1 round=26 year=2015 period=1 form=A subject=civil_law subject_ko=민법 seq=30 number=30
%% source: Q-Net 감정평가사 기출문제 (cst00309) artlSeq=5138001 file=extracted/1008/5138001/제26회_감정평가사 1차_기출/제26회_감정평가사 1차_기출/1교시_감평_A형.pdf
%% conversion: kordoc 4.24.51 -> tools/gen1.py (auto); review=proofread-sonnet; flags=none
\begin{question}{30}
선의취득에 관한 설명으로 옳지 않은 것은? (다툼이 있으면 판례에 따름)
\begin{choices}
  \item 토지는 선의취득의 대상이 되지 못한다.
  \item 점유개정에 의해 간접점유를 취득하였더라도 선의취득을 할 수 없다.
  \item 선의취득자는 임의로 선의취득의 효과를 거부하고 종전 소유자에게 동산을 반환받아 갈 것을 요구할 수 없다.
  \item 점유보조자가 횡령한 동산은 민법 제250조의 도품ㆍ유실물에 해당하지 않는다.
  \item 양수인이 물권적 합의 시점에 선의ㆍ무과실이면, 이후 인도받을 때에 악의이거나 과실이 있더라도 선의취득이 인정된다.
\end{choices}
\end{question}
\answer{5}
